\documentclass[11pt]{article}

\usepackage[margin=1in]{geometry}

\usepackage{booktabs}

\usepackage{adjustbox}

\usepackage{float}

\usepackage{caption}

\captionsetup{font=small,labelfont=bf,skip=4pt}

\title{LLM-as-a-Judge Report: Self-doubt\\\large GSM8K\_20000}

\date{}

\begin{document}

\maketitle

\section*{What was measured?} The judge scores visible self-questioning in the reasoning trace, independently of whether the final answer is correct. A score of 0 means a straight-line trace with no hesitation; 10 means pervasive uncertainty, reversals, or looping. The judge returns the score, a verbatim evidence quote when available, and a one-line summary. The tables below use 0--10 scores and include only valid answered traces from gsm8k\_20000\_judge\_traces.jsonl.

\section*{Examples at different scores} The examples below show the score, final-answer outcome, the judge's evidence, and its one-line summary. An empty evidence field means that the trace contained no suitable verbatim cue for that dimension.

\paragraph{Trace typo50\_real40, index 50.}
\textbf{Score:} 0 on the 0--10 scale. \textbf{Final-answer outcome:} wrong.
\textbf{Judge evidence:} \begin{quote}I don't see any mistakes in my reasoning here.\end{quote}
\textbf{One-line summary:} The trace led to the scores due to a clean straight-line solution with no noticeable self-doubt or misunderstanding of the corrupted words.\medskip
\paragraph{Trace typo25\_real40, index 2.}
\textbf{Score:} 2 on the 0--10 scale. \textbf{Final-answer outcome:} correct.
\textbf{Judge evidence:} \begin{quote}I was initially confused whether the 150\% increase was based on the total cost or the original price\end{quote}
\textbf{One-line summary:} The trace shows a brief hesitation in calculating the increase in value, but overall understands the problem and corrupted words correctly.\medskip
\paragraph{Trace typo50\_real40, index 0.}
\textbf{Score:} 5 on the 0--10 scale. \textbf{Final-answer outcome:} correct.
\textbf{Judge evidence:} \begin{quote}Wait, does that mean she uses four eggs each day to bake muffins? Or is it four eggs per something?\end{quote}
\textbf{One-line summary:} The model's hesitation and questioning of the wording, such as 'Wait, does that mean she uses four eggs each day to bake muffins?' led to the scores.\medskip
\paragraph{Trace typo25\_real70, index 0.}
\textbf{Score:} 8 on the 0--10 scale. \textbf{Final-answer outcome:} wrong.
\textbf{Judge evidence:} \begin{quote}Wait, maybe the problem is that she lays 16 eggs per azy, which is 16 eggs per day.\end{quote}
\textbf{One-line summary:} The trace's repeated reconsideration of the meaning of 'azy' and its potential typos led to high self-doubt and misunderstanding of the question's intended meaning.\medskip

\section*{Aggregate tables}

\begin{table}[H]
\centering
\caption{Mean score by real-word typo ratio}
\small
\begin{adjustbox}{max width=\textwidth}
\begin{tabular}{lrr}
\toprule
\textbf{real\_ratio} & \textbf{n} & \textbf{self\_doubt\_score\_mean} \\
\midrule
10 & 1449 & 5.206 \\
40 & 1474 & 5.288 \\
70 & 1454 & 5.12 \\
\bottomrule
\end{tabular}
\end{adjustbox}
\end{table}

\begin{table}[H]
\centering
\caption{Mean score by typo percentage}
\small
\begin{adjustbox}{max width=\textwidth}
\begin{tabular}{lrr}
\toprule
\textbf{typo\_percentage} & \textbf{n} & \textbf{self\_doubt\_score\_mean} \\
\midrule
0 & 485 & 4.495 \\
25 & 1466 & 4.862 \\
50 & 1472 & 5.12 \\
75 & 1439 & 5.641 \\
\bottomrule
\end{tabular}
\end{adjustbox}
\end{table}

\begin{table}[H]
\centering
\caption{Mean score by final-answer correctness}
\small
\begin{adjustbox}{max width=\textwidth}
\begin{tabular}{lrr}
\toprule
\textbf{outcome} & \textbf{n} & \textbf{self\_doubt\_score\_mean} \\
\midrule
correct & 3816 & 4.691 \\
wrong & 1046 & 6.753 \\
\bottomrule
\end{tabular}
\end{adjustbox}
\end{table}

\section*{Consequences} In this run, self-doubt decreases as the real-word ratio increases: the mean is 1.001 at 10 percent real-word typos, 0.889 at 40 percent, and 0.826 at 70 percent. This is consistent with the interpretation that non-real-word typos are more visibly strange and can trigger more doubt, while real-word typos can be read fluently and silently.

\end{document}